\section{Parte B - Reconstrucción de trayectoria en el plano}

\subsection{Ecuaciones cinemáticas}

La velocidad en el sistema del cuerpo se escribe como $[u_0\ v]^T$. Para expresarla en el sistema inercial se aplica una rotación de ángulo $\psi$:

\begin{equation}
\begin{bmatrix}
\dot{X} \\
\dot{Y}
\end{bmatrix}
=
\begin{bmatrix}
\cos\psi & -\sin\psi \\
\sin\psi & \cos\psi
\end{bmatrix}
\begin{bmatrix}
u_0 \\
v
\end{bmatrix}.
\end{equation}

De esta forma se obtiene

\begin{equation}
\dot{X} = u_0 \cos(\psi) - v \sin(\psi),
\qquad
\dot{Y} = u_0 \sin(\psi) + v \cos(\psi).
\end{equation}

\subsection{Integración de la trayectoria}

La trayectoria se reconstruye a partir de las señales $v(t)$ y $\psi(t)$ obtenidas en la Parte A, con $X(0)=0$ e $Y(0)=0$. La integración numérica se realiza con la misma grilla temporal de simulación.

\begin{table}[H]
\centering
\caption{Condición inicial y posición final reconstruida en la Parte B.}
\label{tab:parte-b-posicion}
\begin{tabular}{lrr}
\toprule
Caso & $X$ [m] & $Y$ [m] \\
\midrule
Inicial & 0.00 & 0.00 \\
Final Parte B & 97.40 & 38.14 \\
\bottomrule
\end{tabular}
\end{table}


\subsection{Trayectoria obtenida}

La Figura~\ref{fig:parte-b-trayectoria} muestra la trayectoria reconstruida en el plano inercial.

\begin{figure}[H]
\centering
\includegraphics[width=0.82\textwidth]{parteB/trayectoria.png}
\caption{Trayectoria reconstruida con el controlador de la Parte A.}
\label{fig:parte-b-trayectoria}
\end{figure}

\subsection{Análisis de la trayectoria}

Se observa en la figura~\ref{fig:parte-b-trayectoria} que en cada cambio de trayectoria, el controlador logra orientar la embarcación correctamente. Sin embargo, como no es posible hacer el ajuste de forma instantánea, la trayectoria real acumula un pequeño error contra la trayectoria perfecta en cada giro. Este comportamiento se puede observar principalmente en el instante $t=40\,s$.

Como se vio en la parte A, con polos situados más a la izquiera se podrían conseguir respuestas más rápidas y reducir el error, a costo de picos altos en la señal de control y cambios más bruscos.

Incluso con un control más fuerte, según las ecuaciones cinemáticas, $v(t)$ afecta a $\dot{X}$ e $\dot{Y}$ con $-v\sin(\psi)$ y $+v\cos(\psi)$, respectivamente. Esto significa que incluso si el ángulo de rumbo $\psi$ fuera exactamente el deseado, la presencia de $v(t)$ hace que la dirección real de desplazamiento difiera levemente de la dirección apuntada por el barco.

\subsection{Discusión}
Si se quisiera controlar directamente la posición en el plano $XY$, sería necesario reformular el problema incorporando realimentación de posición y pasando a MIMO. 
Esto no se podría resolver con el controlador diseñado en la Parte A, ya que la relación entre la entrada de control ($\tau_r$) y las posiciones $X$, $Y$ en el plano es no lineal porque que depende de $\sin(\psi)$ y $\cos(\psi)$.
Para controlar la posición, se tendría que diseñar un controlador no lineal o linealizar el sistema alrededor de un punto para luego aplicar técnicas de control lineal.

\FloatBarrier
